\documentclass[10pt,a4paper]{amsart}
\usepackage[margin=19mm]{geometry}
\usepackage{amsmath,amssymb,mathtools}
\usepackage[scr=rsfs,cal=euler]{mathalfa}
\usepackage{xfrac}
\usepackage[unicode,hidelinks,bookmarks=false]{hyperref}
\newcommand{\ie}{i.e.\ }
\newcommand{\cf}{cf.\ }
\newcommand{\resp}{respectively\ }
\newcommand{\Subsection}{\subsection}
\providecommand{\nobreakdash}{\nobreak\hbox{-}}
\setlength{\emergencystretch}{2em}
\multlinegap0pt
\usepackage{amssymb}
\usepackage{mathtools}
\usepackage{bm}
\usepackage[shortlabels]{enumitem}
\newtheorem{theorem}{Theorem}
\newcommand{\wpF}{\wp}
\title{Painlev\'e Integrability, Hamiltonian Structure\\and Exact Elliptic Reductions of a Generalized Nonlinear Wave Equation}
\author{Alvaro H. Salas}
\date{Source: arXiv:2609.12006v1 (CC BY 4.0)}
\begin{document}
\maketitle

\noindent\emph{Source and licence.} Painlev\'e Integrability, Hamiltonian Structure\\and Exact Elliptic Reductions of a Generalized Nonlinear Wave Equation. Version arXiv:2609.12006v1 (CC BY 4.0), distributed by arXiv under the Creative Commons Attribution 4.0 International licence, http://creativecommons.org/licenses/by/4.0/, as recorded in the arXiv licence metadata for this exact version and shown on the abstract page https://arxiv.org/abs/2609.12006v1. Full licence notice: \texttt{licenses/arxiv-2609.12006-CC-BY-4.0.txt}.\par\smallskip

\noindent\textit{Editorial scope.} Counted unit: the article's theorem eq:quartic (source0.tex lines 331--354). The article's complete original proof of it is delivered with it (source0.tex lines 356--378, one \texttt{proof} environment of 23 lines). No mathematics is written, repaired or re-derived: the statement and the proof are the article's own lines, in the article's own notation, and no retained line is substituted or re-flowed.  No further source-local context is needed: every cross-reference of every retained line resolves inside the two delivered spans.  Nothing was omitted from any retained span, so the package carries no omission record.  No retained line contains a citation, so the wrapper carries no bibliography and no bibliography entry.  No retained line contains a figure environment, an image-inclusion command or drawing code, so no image is delivered and the package declares no asset dependency.  Editorial formatting only, in addition: an AMS \texttt{amsart} preamble that restates the article's own declarations and macros used by the retained lines, namely the environments \texttt{theorem} on separate counters, together with the standard packages those lines need.  The retained environments print in this document as Theorem 1 in order of appearance, whereas the article prints the same items with the numbering of its own full document.  The complete native source of the article as distributed in the arXiv e-print for this exact version is bundled unchanged as \texttt{sources/210067/source0.tex}.
\begin{theorem}[Root-based Weierstrass formula]
Let
\begin{equation}
 (U')^2=R_4(U)
 \label{eq:quartic}
\end{equation}
with $R_4$ a quartic polynomial, and let $U_0$ be a simple root of $R_4$.  Define
\begin{align}
 g_2&=\frac{R_4''(U_0)^2}{48}
      -\frac{R_4'(U_0)R_4'''(U_0)}{24},
 \label{eq:g2}\\
 g_3&=-\frac{R_4''(U_0)^3}{1728}
      +\frac{R_4'(U_0)R_4''(U_0)R_4'''(U_0)}{576}
      -\frac{R_4'(U_0)^2R_4''''(U_0)}{384}.
 \label{eq:g3}
\end{align}
Then every nonconstant local solution represented from that simple root can be written
\begin{equation}
 U(z)=U_0+
 \frac{R_4'(U_0)}{4\left[\wpF(z-z_0;g_2,g_3)-R_4''(U_0)/24\right]},
 \label{eq:weier}
\end{equation}
where $z_0$ is a phase constant.
\end{theorem}

\begin{proof}
Put $U=U_0+1/y$.  Since $R_4(U_0)=0$, multiplying \eqref{eq:quartic} by $y^4$ yields
\begin{equation}
 (y')^2=Ay^3+By^2+Cy+D,
 \label{eq:cubicY}
\end{equation}
where
\[
 A=R_4'(U_0),\quad
 B=\frac{R_4''(U_0)}{2},\quad
 C=\frac{R_4'''(U_0)}{6},\quad
 D=\frac{R_4''''(U_0)}{24}.
\]
The affine change
\[
 y=\frac{4w-B/3}{A}
\]
reduces \eqref{eq:cubicY} to
\[
 (w')^2=4w^3-g_2w-g_3,
\]
with \eqref{eq:g2}--\eqref{eq:g3}.  Hence $w=\wpF(z-z_0;g_2,g_3)$, and solving backward for $U$ gives \eqref{eq:weier}.
\end{proof}

\end{document}
